\section*{الفصل \ref{ch:b1:p-block-16-18} --- الفئة p (2): الأكسجين والهالوجينات والغازات النبيلة}

\begin{solution}{exo:b1:p-block-16-18:1}
\ce{H2S} \babelsublr{$-$II}؛ \ce{S8} 0؛ \ce{SO2} \babelsublr{$+$IV}؛ \ce{SO3^2-} \babelsublr{$+$IV}؛ \ce{SO4^2-}
\babelsublr{$+$VI}؛ \ce{S2O3^2-} \babelsublr{$+$II} في المتوسط؛ \ce{H2S2O7} \babelsublr{$+$VI}.
\end{solution}

\begin{solution}{exo:b1:p-block-16-18:2}
\omterm{def:b1:p-block-16-18:halogen}{الهالوجين} يؤكسد الهاليدات التي تحته: \ce{Cl2 + 2I- -> 2Cl- + I2}
و\ce{Cl2 + 2Br- -> 2Cl- + Br2} يحدثان؛ أما اليود مع البروميد والبروم مع
الكلوريد فلا.
\end{solution}

\begin{solution}{exo:b1:p-block-16-18:3}
\ce{SO2}: الكبريت \omterm{def:b1:lewis-resonance-vsepr:lewis}{بزوج حر}، ورابطة S=O ورابطة \ce{S+-O-} في كل من \omterm{def:b1:lewis-resonance-vsepr:resonance}{بنيتي
الرنين} (أو رابطتان S=O مع ثمانية موسَّعة): منحنٍ.
\ce{SO3}: مثلثي مستوٍ، بثلاث روابط S--O متكافئة. \ce{O3}: ذرة O مركزية
\omterm{def:b1:lewis-resonance-vsepr:lewis}{بزوج حر}، ورابطة O=O ورابطة \babelsublr{O--O$^-$}، \omterm{def:b1:lewis-resonance-vsepr:resonance}{وبنيتا رنين}: منحنٍ.
\end{solution}

\begin{solution}{exo:b1:p-block-16-18:4}
\ce{SO2 + H2O <=> H2SO3}؛ \ce{SO2 + 2NaOH -> Na2SO3 + H2O}؛
\ce{SO3 + H2O -> H2SO4}؛ \ce{SO3 + 2NaOH -> Na2SO4 + H2O}.
\end{solution}

\begin{solution}{exo:b1:p-block-16-18:5}
HF: $x^2/(0.10 - x) = 10^{-3.16}$، $x = \qty{8.0e-3}{mol/L}$، أي pH يساوي 2.10. HCl:
pH يساوي 1.00. فالرابطة H--F قوية وأيون الفلوريد مرتبط بقوة \omterm{def:b1:intermolecular-forces-solvents:hbond}{بروابط
هيدروجينية}: فيكون HF \omterm{def:b1:predominant-reaction:strong-acid}{حمضًا ضعيفًا}.
% ledger: dfg:HF_ao, dfg:F-_ao
\end{solution}

\begin{solution}{exo:b1:p-block-16-18:6}
\ce{BrF5}: خمس روابط \omterm{def:b1:lewis-resonance-vsepr:lewis}{وزوج حر} واحد، هرمي مربع القاعدة. \ce{IF7}:
ثنائي هرم خماسي. \ce{ICl4-}: أربع روابط \omterm{def:b1:lewis-resonance-vsepr:lewis}{وزوجان حران}، مربع
مستوٍ. \ce{XeO3}: ثلاث روابط \omterm{def:b1:lewis-resonance-vsepr:lewis}{وزوج حر} واحد، هرمي مثلثي.
\end{solution}

\begin{solution}{exo:b1:p-block-16-18:7}
$\log K = 2(1.36 - 0.53)/0.059 = 28$. ويكوّن اليود المحرَّر \omterm{def:b1:complexation:complex}{المعقد} الأزرق
الداكن مع النشا.
\end{solution}

\begin{solution}{exo:b1:p-block-16-18:8}
\babelsublr{$E^\circ(\ce{O3}/\ce{O2}) = 2.07 > 0.53$~V}: \ce{O3 + 2I- + 2H+ -> O2 + I2
+ H2O}. ويُعاير اليود المتكوّن بالثيوكبريتات
(\cref{exo:b1:nernst:7}): جزيء \ce{I2} واحد لكل \ce{O3}.
\end{solution}

\begin{solution}{exo:b1:p-block-16-18:9}
\ce{C12H22O11 -> 12C + 11H2O}؛ $M = \qty{342.0}{g/mol}$، $10.0/342.0 =
\qty{0.0292}{mol}$، والكربون $0.0292 \times 12 \times 12.0 = \qty{4.21}{g}$.
\end{solution}

\begin{solution}{exo:b1:p-block-16-18:10}
$x(0.10 + x)/(0.10 - x) = 10^{-1.99}$: $x = \qty{8.6e-3}{mol/L}$،
$[\ce{H3O+}] = 0.1086$، $\mathrm{pH} = 0.96$ بدل 1.00: فالحمضية
الثانية تضيف نحو \qty{9}{\percent} إلى $[\ce{H3O+}]$.
\end{solution}

\begin{solution}{exo:b1:p-block-16-18:11}
تقع المزدوجة $\ce{F2}/\ce{F-}$ (\babelsublr{2.89~V}) فوق كل المزدوجات الأخرى: فلا يستطيع أي \omterm{def:b1:nernst:couple}{مؤكسد} أن يأخذ
الإلكترون من الفلوريد. وفي التحليل الكهربائي المائي، يتفاعل الماء بدلًا منه، إذ يتأكسد عند
جهد أدنى بكثير ($\ce{O2}/\ce{H2O}$، \babelsublr{1.23~V})؛ أما
الفلوريد المنصهر فلا ماء فيه.
\end{solution}

\begin{solution}{exo:b1:p-block-16-18:12}
$\Delta_r G^\circ = -51.4 - (16.40 - 51.57) = \qty{-16.2}{kJ/mol}$،
$\log K = 16.2/5.708 = 2.84$، $K = 7 \times 10^2$.
% ledger: dfg:I2_ao, dfg:I-_ao, dfg:I3-_ao
\end{solution}

\begin{solution}{pb:b1:p-block-16-18:1}
\textbf{1.} 0، \babelsublr{$+$IV}، \babelsublr{$+$VI}، \babelsublr{$+$VI}.
\textbf{2.} \ce{S + O2 -> SO2}.
\textbf{3.} الكبريت بمجالين رابطين \omterm{def:b1:lewis-resonance-vsepr:lewis}{وزوج حر}: منحنٍ، بزاوية نحو
\ang{120}.
\textbf{4.} $10^6/32.1 = \qty{3.12e4}{mol}$.
\textbf{5.} \qty{3.115e4}{mol} من \ce{O2}: $V = 3.115 \times 10^4 \times
8.314 \times 298.15/10^5 = \qty{772}{m^3}$، والهواء $772/0.21 = \qty{3.7e3}{m^3}$.
\textbf{6.} الاحتراق ناشر للحرارة جدًا؛ \omterm{def:b1:elementary-steps:catalyst}{والحفّاز} يعمل في مجال
محدود من درجات الحرارة، وأكسدة \ce{SO2}، الناشرة للحرارة أيضًا، تكون أقل
اكتمالًا على الساخن.
\textbf{7.} \ce{2SO2 + O2 <=> 2SO3}.
\textbf{8.} ثابت التوازن الكبير لا يقول شيئًا عن السرعة: فعند
درجة حرارة الغرفة لا يجري التفاعل.
\textbf{9.} يوفّر مسارًا \omterm{def:b1:rate-laws:activation-energy}{بطاقة تنشيط} أدنى ويُستعاد:
فهو \omterm{def:b1:elementary-steps:catalyst}{حفّاز} غير متجانس.
\textbf{10.} مثلثي مستوٍ.
\textbf{11.} مزيد من الأكسجين يُبقي حاصل التفاعل
$Q = p_{\ce{SO3}}^2/(p_{\ce{SO2}}^2 p_{\ce{O2}})$ أدنى: فيمضي تحول
\ce{SO2} أبعد.
\textbf{12.} \qty{0.3}{\percent}.
\textbf{13.} \ce{SO3 + H2SO4 -> H2S2O7}.
\textbf{14.} يكوّن ضبابًا من قطيرات الحمض الدقيقة يمر عبر
جهاز الامتصاص.
\textbf{15.} \ce{H2S2O7 + H2O -> 2H2SO4}.
\textbf{16.} \ce{2S + 3O2 + 2H2O -> 2H2SO4}.
\textbf{17.} جزيء ماء واحد لكل ذرة كبريت: $\num{3.115e4} \times 18.0 =
\qty{561}{kg}$.
\textbf{18.} يطلق التخفيف حرارة كبيرة؛ والماء المصبوب على الحمض
يغلي فورًا ويرشّ الحمض.
\textbf{19.} $x(0.050 + x)/(0.050 - x) = 10^{-1.99}$: $x = \qty{7.5e-3}{mol/L}$،
$[\ce{H3O+}] = \qty{0.0575}{mol/L}$، $\mathrm{pH} = 1.24$.
\textbf{20.} $3.12 \times 10^4 \times 0.995 \times 98.1 = \qty{3.04}{t}$.
\textbf{21.} $84 \times 98.1/32.1 = \qty{257}{Mt}$.
\textbf{22.} $98.1/32.1 =$ \textbf{\qty{3.06}{t} من حمض الكبريتيك لكل طن من
الكبريت}.
\end{solution}
